\section{Constructing Proofs}

Although the rules are pretty simple, building proofs can get quite complicated.
Much like chess, the rules are easy to learn, but yon can get into trouble
quickly if you don't adopt some basic straegies to guide your moves.

We can tersely describe the basic strategy for TFL proofs as follows.

\begin{outline}[enumerate]

  \1 Work backwards from the conclusion.
  \2 Conjunction: what components needed to use the \eand I rule as the
  last step? i.e. what are the conjuncts?
  \2 Conditional: what are the assumption and conclusion of the subproof
  we need to take the last step? i.e. what are the antecedant and consequent
  of the conditional?
  \2 Negation: what subproof is needed to show a contradiction?
  \2 Disjunction: which of the disjuncts looks easiest to prove?

  \1 Work forwards from the current step.
  \2 Conjunction: unpack the conjuncts
  \2 Disjunction: what sentence do you need to prove from both disjuncts?
  \2 Conditional: how can you prove the antecedant?
  \2 Negation: how can you prove a contradiction?

  \1 Leverage \ebot
  \2 Work toward a contradiction to conclude anything, especially in a subproof
  \2 Indirect proof: build a contradictory subproof that assumes the negation

\end{outline}

Ideally, we will apply these stragegies in this order, wherever the rules are
appropriate to the current state of the proof. The \emph{work backwards} rules
in particular are intended as scaffolding to indicate which portions of the
overall proof need to be filled in. The rules themselves reflect this
partially-completed aspect.

\Needspace{7\baselineskip}
\begin{fitchrule}[Proving Conjunctions]
  \label{def:proving-conjunctions}
  \leavevmode
  \vspace{-0.9\baselineskip}
  \begin{fitchproof}
    \have[ ]     {1} {\vdots}    
    \have[m]     {2} {\ephi}
    \have[ ]     {3} {\vdots}    
    \have[n]     {4} {\epsi}   
    \have[n][+1] {5} {\ephi\eand\epsi} \ai{2,4}
  \end{fitchproof}
\end{fitchrule}

\Needspace{5\baselineskip}
\twoitems{

\begin{fitchrule}[Proving Conditionals]
  \label{def:proving-conditionals}
  \leavevmode
  \vspace{-0.9\baselineskip}
  \begin{fitchproof}
    \open
    \hypo[m]     {1} {\ephi}          \by{AS}{}
    \have[ ]     {2} {\vdots}
    \have[n]     {3} {\epsi}
    \close
    \have[n][+1] {4} {\ephi\eif\epsi} \ii{1-3}
  \end{fitchproof}
\end{fitchrule}

}{
  
\begin{fitchrule}[Proving Negations]
  \label{def:proving-negations}
  \leavevmode
  \vspace{-0.9\baselineskip}
  \begin{fitchproof}
    \open
    \hypo[m]     {1} {\ephi}      \by{AS}{}
    \have[ ]     {2} {\vdots}
    \have[n]     {3} {\ebot}
    \close
    \have[n][+1] {4} {\enot\ephi} \ni{1-3}
  \end{fitchproof}  
\end{fitchrule}

}

Working backward from disjunctions is possible too, although the ambiguity
may lead to some backtracking. So this is a last resort.

\Needspace{4\baselineskip}
\begin{fitchrule}[Proving Disjunctions] There are two possible justifications
  for a disjunction, i.e. each of its disjuncts:\\
  \label{def:proving-disjunctionss}
  \leavevmode
  %\vspace{-0.9\baselineskip}
  \twoitems{
    \begin{fitchproof}
      \have[ ]     {1} {\vdots}    
      \have[n]     {2} {\ephi}
      \have[n][+1] {3} {\ephi\eor\epsi} \oi{2}
    \end{fitchproof}
  }{
    \begin{fitchproof}
      \have[ ]     {1} {\vdots}    
      \have[n]     {2} {\epsi}
      \have[n][+1] {3} {\ephi\eor\epsi} \oi{2}
    \end{fitchproof}
  }
\end{fitchrule}

Working forward can be guided by the same schematic rules. For conjunction we
just unpack both of the conjuncts.

\begin{fitchrule}[Unpacking Conjunctions]
  \label{def:unpacking-conjunctions}
  \leavevmode
  \vspace{-0.9\baselineskip}
  \begin{fitchproof}
    \have[n]     {1} {\ephi\eand\epsi} 
    \have[ ][+1] {2} {\ephi}
    \have[n][+2] {3} {\epsi}
  \end{fitchproof}  
\end{fitchrule}

For disjunction we need to think more about what we are trying to do. Suppose
that \ephi and \epsi are the disjuncts and \echi is what we are trying to
prove (or a step in that direction).

\Needspace{10\baselineskip}
\begin{fitchrule}[Unpacking Disjunctions]
  \label{def:unpacking-disjunctions}
  \leavevmode
  \vspace{-0.9\baselineskip}
  \begin{fitchproof}
    \have[k]     {1} {\ephi\eor\epsi}
    \open
    \hypo[k][+1] {2} {\ephi} \by{AS}{}
    \have[ ]     {3} {\vdots}
    \have[m]     {4} {\echi}
    \close    
    \open
    \hypo[m][+1] {5} {\ephi} \by{AS}{}
    \have[ ]     {6} {\vdots}
    \have[n]     {7} {\echi}
    \close
    \have[n][+1] {8} {\echi} \oe{1,2-4,5-7}
  \end{fitchproof}  
\end{fitchrule}

Finally, conditionals and negations are pretty easy to unpack.

\Needspace{6\baselineskip}
\twoitems{

\begin{fitchrule}[Unpacking Conditionals]
  \label{def:unpacking-conditionals}
  \leavevmode
  \vspace{-0.9\baselineskip}
  \begin{fitchproof}
    \have[m]     {1} {\ephi\eif\epsi} 
    \have[ ]     {2} {\vdots}
    \have[n]     {3} {\ephi}
    \have[n][+1] {4} {\epsi} \ie{1,3}
  \end{fitchproof}    
\end{fitchrule}

} {
  
\begin{fitchrule}[Unpacking Negations]
  \label{def:unpacking-negationss}
  \leavevmode
  \vspace{-0.9\baselineskip}
  \begin{fitchproof}
    \have[m]     {1} {\enot\ephi} 
    \have[ ]     {2} {\vdots}
    \have[n]     {3} {\ephi}
    \have[n][+1] {4} {\ebot} \ne{1,3}
  \end{fitchproof}      
\end{fitchrule}

}

Sometimes the best way to get a subproof to its intended conclusion is to work
toward a contradiction. Then you can conclude the subproof with literally
anything, which is a powerfully flexible tool. Consider this example from
\cite{Mag25}[p156] which shows how finding contradictions can be helpful,
especially when working with disjunctions.

\begin{fitchproof}
  \hypo {1} {A \eor B}  \by{PR}{} 
  \hypo {2} {\enot A}   \by{PR}{}
  \open
  \hypo {3} {A}
  \have {4} {\ebot}     \ne{2,3}
  \have {5} {B}         \by{X}{4}
  \close
  \open
  \hypo {6} {B}
  \have {7} {B}         \by{R}{6}
  \close
  \have {7} {B}         \by{R}{1,3-5,6-7}
\end{fitchproof}      

Similarly, we can get any sentence we want in the main proof (or current
subproof) by assuming its negation and proving a contradiction. We will call
this technique \emph{indirect proof}.

\begin{fitchproof}
  \open
  \hypo {1} {\enot(A\eor\enot A)} \by{AS}{}
  \open
  \hypo {2} {A}                   \by{AS}{}
  \have {3} {A\eor\enot A}        \oi{2}
  \have {4} {\ebot}               \ne{1,3}
  \close
  \have {5} {\enot A}             \ni{2-4}
  \have {6} {A \eor\enot A}       \oi{5}
  \have {7} {\ebot}               \ne{1,6}
  \close
  \have {8} {A \eor\enot A}        \by{IP}{1-7}

\end{fitchproof}      

Later we will use these last two valid arguments, among others, as the basis of
\emph{derived rules} that will let us skip many of these tedious steps when
we write a proof that needs them. Notice how the second one--the
\emph{excluded middle} that we saw last week--does not require any premises at
all. The conclusion $A\eor\not A$ is a tautology.

For now, we can formulate a rule for indirect proof (IP) as it was demonstrated
in the last example.

\begin{fitchrule}[Indirect Proof]
  \label{def:indirect-proof-rule}
  \leavevmode
  \vspace{-0.9\baselineskip}
  \begin{fitchproof}
    \open
    \hypo[m]     {1} {\enot\ephi} \by{AS}{}
    \have[ ]     {2} {\vdots}
    \have[n]     {3} {\ebot}
    \close
    \have[n][+1] {4} {\phi}       \by{IP}{1-3}
  \end{fitchproof}      
\end{fitchrule}

